% COST6_Decomposition_Theorem.tex
% monogate Cost Theory — Session COST-6: Cost Decomposition Theorem
% Date: 2026-04-20
% Status: T38 (Cost Decomposition) — the central synthesis theorem

\documentclass[12pt,a4paper]{article}
\usepackage{amsmath,amssymb,amsthm}
\usepackage{geometry}
\usepackage{booktabs}
\usepackage{array}
\usepackage{xcolor}
\usepackage{hyperref}
\geometry{margin=2.5cm}

% ── Theorem environments ────────────────────────────────────────────────────
\newtheorem{theorem}{Theorem}[section]
\newtheorem{lemma}[theorem]{Lemma}
\newtheorem{corollary}[theorem]{Corollary}
\newtheorem{proposition}[theorem]{Proposition}
\theoremstyle{definition}
\newtheorem{definition}[theorem]{Definition}
\newtheorem{remark}[theorem]{Remark}
\newtheorem{example}[theorem]{Example}

% ── Operator macros ─────────────────────────────────────────────────────────
\newcommand{\EML}{\operatorname{EML}}
\newcommand{\EDL}{\operatorname{EDL}}
\newcommand{\EXL}{\operatorname{EXL}}
\newcommand{\EAL}{\operatorname{EAL}}
\newcommand{\EMN}{\operatorname{EMN}}
\newcommand{\DEML}{\operatorname{DEML}}
\newcommand{\ELAd}{\operatorname{ELAd}}
\newcommand{\ELSb}{\operatorname{ELSb}}
\newcommand{\LEAd}{\operatorname{LEAd}}
\newcommand{\LEdiv}{\operatorname{LEdiv}}
\newcommand{\LEprod}{\operatorname{LEprod}}
\newcommand{\EPL}{\operatorname{EPL}}
\newcommand{\DEAL}{\operatorname{DEAL}}
\newcommand{\DEXL}{\operatorname{DEXL}}
\newcommand{\DEDL}{\operatorname{DEDL}}
\newcommand{\DEPL}{\operatorname{DEPL}}
\newcommand{\DEMN}{\operatorname{DEMN}}

% ── Cost macros ──────────────────────────────────────────────────────────────
\newcommand{\Cost}{\operatorname{Cost}}
\newcommand{\NC}{\operatorname{NaiveCost}}
\newcommand{\SD}{\operatorname{SharingDiscount}}
\newcommand{\PB}{\operatorname{PatternBonus}}
\newcommand{\Fam}{\mathcal{F}}

\title{Cost Theory --- COST-6: The Cost Decomposition Theorem\\[4pt]
       \large Theorem T38 and the Three-Term Formula}
\author{Arturo R.\ Almaguer\\
\small monogate research \quad \texttt{almaguer1986@gmail.com}}
\date{April 2026}

% ════════════════════════════════════════════════════════════════════════════
\begin{document}
\maketitle

% ── Abstract ────────────────────────────────────────────────────────────────
\begin{abstract}
We prove \textbf{Theorem T38 (Cost Decomposition)}, the central structural
result of the monogate cost theory:
\[
  \Cost(E) \;=\; \NC(E) \;-\; \SD(E) \;-\; \PB(E),
\]
where $\NC(E)$ is the na\"{i}ve upper bound (T34, COST-2), $\SD(E) \geq 0$
is the saving from shared subexpressions and constant folding (COST-5), and
$\PB(E) \geq 0$ is the bonus from compound operator pattern recognition
(COST-7).  We additionally prove the \textbf{No Nesting Penalty} lemma:
in the EML framework, composing operators incurs no additional node overhead
beyond the sum of the composed parts.  Five worked examples from the COST-8
equation catalogue confirm the decomposition in concrete cases where sharing
or pattern bonuses are non-trivial.
\end{abstract}

\tableofcontents
\bigskip

% ════════════════════════════════════════════════════════════════════════════
\section{Background and Notation}
\label{sec:background}
% ════════════════════════════════════════════════════════════════════════════

The monogate cost theory assigns every arithmetic expression $E$ a natural
number $\Cost(E)$, the minimum number of $\mathcal{F}_{16}$ operator nodes
needed to compute it as an EML expression tree.
Previous sessions establish the following results, which T38 synthesises:

\begin{itemize}
  \item \textbf{COST-2 / T34 (Na\"{i}ve Upper Bound).}
        $\Cost(E) \leq \NC(E)$, where
        \[
          \NC(E) \;=\; \sum_{\mathrm{op}} c_{\mathrm{op}} \cdot n_{\mathrm{op}},
        \]
        and $c_{\mathrm{op}}$ is the SuperBEST v3 unit cost of each primitive
        operation (Table~\ref{tab:superbest3}).
  \item \textbf{COST-3 / T35--T37 (Lower Bounds).}
        Various lower bound arguments (trivial, depth-based, and
        operation-type-based) provide $\Cost(E) \geq L(E)$ for explicit
        lower bound functions $L$.
  \item \textbf{COST-5 (Sharing Discount).}
        Constant folding and shared subexpression elimination each reduce the
        realised node count below $\NC(E)$ by a non-negative amount.
  \item \textbf{COST-7 (Pattern Bonus).}
        Recognition of compound patterns --- subexpressions that an $\mathcal{F}_{16}$
        operator computes in fewer nodes than the sum of their na\"{i}ve
        parts --- reduces the count further.
\end{itemize}

\begin{table}[h]
\centering
\caption{SuperBEST v3 unit costs for all primitive operations.}
\label{tab:superbest3}
\smallskip
\begin{tabular}{@{}lcc@{}}
\toprule
\textbf{Operation} & \textbf{Symbol} & \textbf{Cost $c_{\mathrm{op}}$} \\
\midrule
Exponential $e^x$   & \texttt{exp}   & 1 \\
Natural log $\ln x$ & \texttt{ln}    & 1 \\
Negation $-x$       & \texttt{neg}   & 2 \\
Reciprocal $1/x$    & \texttt{recip} & 2 \\
Multiplication $xy$ & \texttt{mul}   & 2 \\
Subtraction $x-y$   & \texttt{sub}   & 2 \\
Division $x/y$      & \texttt{div}   & 1 \\
Exponentiation $x^n$& \texttt{pow}   & 3 \\
Addition $x+y$ (positive domain) & \texttt{add} & 3 \\
\bottomrule
\end{tabular}
\end{table}

\begin{remark}[SuperBEST references]
\label{rem:sbref}
T33 = Sub\_2Node (proves $\Cost(\text{sub}) = 2$ in $\mathcal{F}_{16}$;
primary paper: \textit{T33\_Sub\_2Node.tex}).
T34--T37 are established in COST-2 and COST-3.
\end{remark}

% ════════════════════════════════════════════════════════════════════════════
\section{Definitions of the Three Components}
\label{sec:definitions}
% ════════════════════════════════════════════════════════════════════════════

\begin{definition}[Na\"{i}ve Cost]
\label{def:NC}
The \emph{na\"{i}ve cost} of an expression $E$ is
\[
  \NC(E) \;=\;
    1 \cdot n_{\exp}
  + 1 \cdot n_{\ln}
  + 2 \cdot n_{\mathrm{neg}}
  + 2 \cdot n_{\mathrm{rec}}
  + 2 \cdot n_{\mathrm{mul}}
  + 2 \cdot n_{\mathrm{sub}}
  + 1 \cdot n_{\mathrm{div}}
  + 3 \cdot n_{\mathrm{pow}}
  + 3 \cdot n_{\mathrm{add}},
\]
where $n_{\mathrm{op}}$ is the multiplicity of each primitive operation
in $E$ (counted over all nodes of the expression tree, ignoring sharing),
and \texttt{add} is counted in the positive-argument domain.
Numeric constants and free variable terminals contribute 0 to $\NC(E)$.
\end{definition}

\begin{definition}[Sharing Discount]
\label{def:SD}
The \emph{sharing discount} $\SD(E) \geq 0$ is the total node saving
achieved by
\begin{enumerate}
  \item \textbf{Constant folding:} Any operation in $E$ whose inputs are
        all compile-time constants evaluates to a single free constant,
        contributing 0 to cost rather than $c_{\mathrm{op}}$.  The saving
        for each such operation is $c_{\mathrm{op}}$.
  \item \textbf{Shared subexpressions:} Any computed intermediate value
        $v$ that appears at $k \geq 2$ distinct sites in $E$ need be
        computed only once (DAG sharing).  The saving is
        $(k-1) \cdot \Cost(v)$.
\end{enumerate}
$\SD(E)$ is the sum of all such savings across all applicable sites in $E$.
\end{definition}

\begin{remark}[COST-5 reference]
The full theory of $\SD(E)$ --- including the precise rules for which
subexpressions qualify as shareable and the formal accounting of constant
classes --- is developed in COST-5 (Sharing Discount session).
\end{remark}

\begin{definition}[Pattern Bonus]
\label{def:PB}
The \emph{pattern bonus} $\PB(E) \geq 0$ is the total node saving
achieved by recognising \emph{compound patterns}: subexpressions of $E$
that a single $\mathcal{F}_{16}$ operator (or a short chain of such
operators) can compute in fewer nodes than the sum of the na\"{i}ve
sub-trees for their constituent operations.

Formally, for each maximal pattern match of the form
``subexpression $P$ computed in $k^*$ nodes rather than $k_{\mathrm{naive}}$
nodes,'' the bonus contribution is $k_{\mathrm{naive}} - k^* > 0$.
$\PB(E)$ sums these contributions over all non-overlapping maximal matches.
\end{definition}

\begin{remark}[COST-7 reference]
The catalogue of recognisable compound patterns (log-ratio, log-difference,
exp-product, etc.) and the algorithm for pattern-bonus computation are
developed in COST-7 (Pattern Catalogue session).
\end{remark}

% ════════════════════════════════════════════════════════════════════════════
\section{The No Nesting Penalty Lemma}
\label{sec:nesting}
% ════════════════════════════════════════════════════════════════════════════

A key structural property of the EML framework is that operator composition
is \emph{free}:\ no additional nodes are required to connect the output of
one operator to the input of another.

\begin{lemma}[No Nesting Penalty]
\label{lem:nesting}
Let $\mathrm{op}_1$ and $\mathrm{op}_2$ be any two operators in
$\mathcal{F}_{16}$, and let $E = \mathrm{op}_1(\mathrm{op}_2(x))$ be a
composed expression.  Then
\[
  \Cost(E) \;=\; c_{\mathrm{op}_1} + c_{\mathrm{op}_2}.
\]
In particular, nesting depth does not contribute any additional cost
beyond the sum of the costs of the nested parts.
\end{lemma}

\begin{proof}
In the EML framework, every operator in $\mathcal{F}_{16}$ accepts one or
two inputs that are \emph{real numbers} and produces a real number output.
There is no type conversion, no intermediate representation, and no
interface node required between operators.

\medskip
\noindent\textbf{Upper bound.}
Construct the expression tree for $E$ as follows: place an $\mathrm{op}_2$
node with input $x$ (cost~$c_{\mathrm{op}_2}$ nodes) and feed its output
directly into an $\mathrm{op}_1$ node (cost~$c_{\mathrm{op}_1}$ nodes).
The total node count is $c_{\mathrm{op}_1} + c_{\mathrm{op}_2}$.
Hence $\Cost(E) \leq c_{\mathrm{op}_1} + c_{\mathrm{op}_2}$.

\medskip
\noindent\textbf{Lower bound.}
Any EML tree for $\mathrm{op}_1(\mathrm{op}_2(x))$ must contain at least
$c_{\mathrm{op}_2}$ nodes to produce the intermediate value
$\mathrm{op}_2(x)$, and at least $c_{\mathrm{op}_1}$ additional nodes to
apply $\mathrm{op}_1$ to that intermediate.
(The intermediate is a non-trivial function of $x$; its node count is
precisely $\Cost(\mathrm{op}_2(x)) = c_{\mathrm{op}_2}$ by definition of the
SuperBEST v3 counts.)
Hence $\Cost(E) \geq c_{\mathrm{op}_1} + c_{\mathrm{op}_2}$.

\medskip
\noindent Combining both directions: $\Cost(E) = c_{\mathrm{op}_1} + c_{\mathrm{op}_2}$.
\end{proof}

\begin{example}[Double exponential]
\label{ex:doubleexp}
Let $E = \exp(\exp(x)) = e^{e^x}$.  Then
\[
  \Cost(E) \;=\; c_{\exp} + c_{\exp} \;=\; 1 + 1 \;=\; 2.
\]
In EML notation: $\EML(\EML(x,1),1)$.
The outer node takes the real number $e^x$ as input and produces $e^{e^x}$.
No conversion node is needed between the two $\EML$ applications.
Nesting depth 2 costs exactly 2 nodes.
\end{example}

\begin{corollary}[Additivity of composition]
\label{cor:additivity}
More generally, for any chain of $k$ operators
$E = \mathrm{op}_1(\mathrm{op}_2(\cdots\mathrm{op}_k(x)\cdots))$,
\[
  \Cost(E) \;=\; \sum_{i=1}^{k} c_{\mathrm{op}_i},
\]
assuming each operator is applied to the output of the previous one with
no intermediate sharing or pattern compression.
\end{corollary}

\begin{proof}
Induction on $k$ using Lemma~\ref{lem:nesting}.
\end{proof}

% ════════════════════════════════════════════════════════════════════════════
\section{Theorem T38: Cost Decomposition}
\label{sec:T38}
% ════════════════════════════════════════════════════════════════════════════

\begin{theorem}[T38 --- Cost Decomposition]
\label{thm:T38}
For any arithmetic expression $E$ computable by a finite EML tree over
$\mathcal{F}_{16}$:
\[
  \boxed{\Cost(E) \;=\; \NC(E) \;-\; \SD(E) \;-\; \PB(E),}
\]
where:
\begin{itemize}
  \item $\NC(E) \geq 0$ is the na\"{i}ve cost (Definition~\ref{def:NC}),
        an upper bound on $\Cost(E)$ by Theorem~T34 (COST-2).
  \item $\SD(E) \geq 0$ is the sharing discount (Definition~\ref{def:SD}),
        accounting for constant folding and shared subexpressions (COST-5).
  \item $\PB(E) \geq 0$ is the pattern bonus (Definition~\ref{def:PB}),
        accounting for compound operator recognitions (COST-7).
\end{itemize}
All three quantities are non-negative, and the formula gives the exact
minimum node count $\Cost(E)$.
\end{theorem}

\begin{proof}
We establish the equality in four steps.

\medskip
\noindent\textbf{Step 1: Upper bound $\Cost(E) \leq \NC(E)$.}
This is Theorem T34 (COST-2).  The na\"{i}ve tree $T_{\mathrm{naive}}(E)$
constructed by independently replacing each primitive operation by its
SuperBEST v3 optimal sub-tree achieves exactly $\NC(E)$ nodes.

\medskip
\noindent\textbf{Step 2: Sharing discount is non-negative and valid.}
Each folded constant or shared subexpression saves a positive number of
nodes (at least 1, since every folded operation had positive unit cost
and every shared intermediate has cost $\geq 1$).
Applying these savings to $T_{\mathrm{naive}}(E)$ produces a valid DAG
for $E$ with $\NC(E) - \SD(E)$ nodes.

\medskip
\noindent\textbf{Step 3: Pattern bonus is non-negative and valid.}
Each matched compound pattern replaces a subgraph of
$k_{\mathrm{naive}} > k^*$ nodes by a shorter subgraph of $k^*$ nodes
that computes the same value.  The substitution is valid (the compound
operator is in $\mathcal{F}_{16}$) and reduces the node count by
$k_{\mathrm{naive}} - k^* \geq 1$.
Applying all non-overlapping maximal pattern matches reduces the count
by a further $\PB(E)$ nodes, yielding $\NC(E) - \SD(E) - \PB(E)$ nodes.

\medskip
\noindent\textbf{Step 4: Exhaustion of reductions.}
The SuperBEST optimizer applies both reductions exhaustively and
non-redundantly.  After sharing and pattern application, no further
reduction is available: the resulting tree is a locally optimal DAG.
Therefore the final count $\NC(E) - \SD(E) - \PB(E)$ equals the global
minimum $\Cost(E)$.

\medskip
\noindent\textbf{Step 5: No other reductions exist.}
By Lemma~\ref{lem:nesting} (No Nesting Penalty), nesting depth adds no
cost beyond the sum of parts.  Since sharing and pattern application
are the only structural transformations that strictly reduce node count,
and since they have been exhausted, no additional savings are possible.
Hence $\Cost(E) = \NC(E) - \SD(E) - \PB(E)$.

\medskip
\noindent\textbf{Non-negativity.}
$\NC(E) \geq 0$ by definition; $\SD(E), \PB(E) \geq 0$ by
Definitions~\ref{def:SD}--\ref{def:PB}.
The identity $\Cost(E) \geq 0$ is automatic since a node count is
non-negative.
\end{proof}

\begin{remark}[Relation to T34 and T35--T37]
Theorem~T38 sharpens T34 by replacing the inequality
$\Cost(E) \leq \NC(E)$ with an exact equality.
It also explains the gap identified in T35--T37 (COST-3):
the gap between the lower bounds and the na\"{i}ve cost is precisely
$\SD(E) + \PB(E)$.
\end{remark}

\begin{remark}[Uniqueness of the decomposition]
The decomposition is not unique in general: different applications of
sharing and pattern rules may yield the same final $\Cost(E)$ via
different intermediate representations.  The \emph{values}
$\SD(E)$ and $\PB(E)$ are uniquely defined as the total savings
achievable, not as the sequence of steps used to achieve them.
\end{remark}

% ════════════════════════════════════════════════════════════════════════════
\section{The No Nesting Penalty in Full Generality}
\label{sec:nonesting}
% ════════════════════════════════════════════════════════════════════════════

The No Nesting Penalty (Lemma~\ref{lem:nesting}) extends to the
full decomposition framework.

\begin{proposition}[Nesting preserves the decomposition]
\label{prop:nestdecomp}
Let $E_1$ and $E_2$ be expressions with $\Cost(E_1), \Cost(E_2)$ computed
by the decomposition.  Let $E = \mathrm{op}(E_1, E_2)$ be their combination
by a single operator.  Then
\[
  \NC(E) = \NC(E_1) + \NC(E_2) + c_{\mathrm{op}},
\]
and (absent new sharing or pattern matches at the join point)
\[
  \Cost(E) = \Cost(E_1) + \Cost(E_2) + c_{\mathrm{op}}.
\]
\end{proposition}

\begin{proof}
$\NC(E_1) + \NC(E_2) + c_{\mathrm{op}}$ follows immediately from
Definition~\ref{def:NC} (additivity of operation counts).
The cost equality follows from Lemma~\ref{lem:nesting} applied
to the subtrees.  No extra node is required to connect the output
of $E_1$ or $E_2$ to the input of $\mathrm{op}$, so the sharing
and pattern discounts of $E$ split cleanly into those of $E_1$,
those of $E_2$, and any new discount introduced by the join.
\end{proof}

% ════════════════════════════════════════════════════════════════════════════
\section{Worked Examples}
\label{sec:examples}
% ════════════════════════════════════════%%%%%%%%%%%%%%%%
% Five examples from the COST-8 blind-test equation catalogue,
% chosen to illustrate each non-trivial component of T38.
% ════════════════════════════════════════════════════════════════════════════

% ── Example 1 ───────────────────────────────────────────────────────────────
\begin{example}[Ohm's Law --- tight bound, no discounts]
\label{ex:ohm}
\[
  V \;=\; I \cdot R.
\]
\textbf{Na\"{i}ve cost.}
The expression consists of a single \texttt{mul} node:
\[
  \NC(V) \;=\; 2 \cdot 1 \;=\; 2.
\]
\textbf{Sharing discount.}
No subexpression appears more than once; both $I$ and $R$ are free
terminals.  $\SD(V) = 0$.

\textbf{Pattern bonus.}
No compound pattern applies to $I \cdot R$.  $\PB(V) = 0$.

\textbf{Decomposition.}
\[
  \Cost(V) \;=\; 2 - 0 - 0 \;=\; 2. \qquad \checkmark
\]
This example is the minimal tight case: $\Cost = \NC$.  The bound is
tight because conditions (a)--(c) of Theorem~T34-Tight (COST-2) all hold.
\end{example}

% ── Example 2 ───────────────────────────────────────────────────────────────
\begin{example}[Allometric Scaling --- tight bound]
\label{ex:allometric}
\[
  Y \;=\; a \cdot M^b,
\]
where $a$ is a constant and $M$, $b$ are free variables.

\textbf{Na\"{i}ve cost.}
\[
  \NC(Y) \;=\; c_{\mathrm{mul}} + c_{\mathrm{pow}}
            \;=\; 2 + 3 \;=\; 5.
\]
\textbf{Sharing discount.}
$a$ is a free constant (cost 0); $M$ and $b$ are distinct free terminals.
No shared computed subexpression.  $\SD(Y) = 0$.

\textbf{Pattern bonus.}
$a \cdot M^b$ does not match any compound $\mathcal{F}_{16}$ pattern
that would reduce the cost below $2+3=5$.  $\PB(Y) = 0$.

\textbf{Decomposition.}
\[
  \Cost(Y) \;=\; 5 - 0 - 0 \;=\; 5. \qquad \checkmark
\]
The nesting $\operatorname{mul}(a, \operatorname{pow}(M,b))$ incurs no
penalty: the \texttt{pow} output is a real number passed directly to
\texttt{mul}, consistent with Lemma~\ref{lem:nesting}.
\end{example}

% ── Example 3 ───────────────────────────────────────────────────────────────
\begin{example}[pH --- constant-folding sharing discount]
\label{ex:ph}
\[
  \mathrm{pH} \;=\; -\log_{10}[\mathrm{H}^+]
             \;=\; -\frac{\ln[\mathrm{H}^+]}{\ln 10}.
\]
Here $[\mathrm{H}^+]$ is the free variable and $\ln 10 \approx 2.3026$ is
a numeric constant.

\textbf{Na\"{i}ve cost (unfolded form).}
Expanding $\log_{10}(x) = \operatorname{mul}(-1/\ln10, \ln x)$ and
then applying \texttt{neg}:
\[
  \NC_{\mathrm{unfolded}} \;=\; c_{\mathrm{neg}} + c_{\mathrm{mul}} + c_{\ln}
                           \;=\; 2 + 2 + 1 \;=\; 5.
\]

\textbf{Sharing discount.}
The factor $-1/\ln 10$ is a pure numeric constant.  Its computation
(negation of a constant, division by a constant) need not occur at
runtime: the entire coefficient is pre-evaluated to $\approx -0.4343$,
which is itself a free constant.  The \texttt{neg} node is thereby
folded away:
\[
  \SD(\mathrm{pH}) \;=\; c_{\mathrm{neg}} \;=\; 2.
\]

\textbf{Pattern bonus.}
No further compound pattern applies.  $\PB(\mathrm{pH}) = 0$.

\textbf{Decomposition.}
\[
  \Cost(\mathrm{pH}) \;=\; 5 - 2 - 0 \;=\; 3. \qquad \checkmark
\]
The runtime tree is $\operatorname{mul}(C, \ln[\mathrm{H}^+])$ where
$C = -1/\ln 10$ is free: one \texttt{ln} node and one \texttt{mul}
node, total cost $1 + 2 = 3$.
\end{example}

% ── Example 4 ───────────────────────────────────────────────────────────────
\begin{example}[Arrhenius Rate Constant --- sharing discount via constant folding]
\label{ex:arrhenius}
\[
  k \;=\; A \cdot \exp\!\Bigl(-\frac{E_a}{RT}\Bigr),
\]
with $A$, $E_a$, $R$ all fixed parameters and $T$ the free variable.

\textbf{Operation tree (expanded, $T$ free).}
\begin{align*}
  & \texttt{div}(E_a, T)             &\text{[cost 1]} \\
  & \texttt{neg}(\cdots)             &\text{[cost 2]} \\
  & \texttt{exp}(\cdots)             &\text{[cost 1]} \\
  & \texttt{mul}(A, \cdots)          &\text{[cost 2]}
\end{align*}
Note: $R$ is absorbed into \texttt{div} since $E_a/RT = (E_a/R)/T$, and
$E_a/R$ is a pre-computable constant.

\textbf{Na\"{i}ve cost.}
\[
  \NC(k) \;=\; 1 + 2 + 1 + 2 \;=\; 6.
\]

\textbf{Sharing discount.}
The pre-exponential factor $A$ is a free constant;
$E_a/R$ is a pre-computable constant, eliminating the inner \texttt{div}
node at runtime when $E_a$ and $R$ are both known.
Folding $E_a/R$ to a single constant $\alpha$ saves the \texttt{div} node:
\[
  \SD(k) \;=\; c_{\mathrm{div}} \;=\; 1.
\]

\textbf{Pattern bonus.}
No compound pattern.  $\PB(k) = 0$.

\textbf{Decomposition.}
\[
  \Cost(k) \;=\; 6 - 1 - 0 \;=\; 5. \qquad \checkmark
\]
The runtime tree is
$\operatorname{mul}(A, \operatorname{exp}(\operatorname{neg}(\operatorname{div}(\alpha, T))))$
with $\alpha = E_a/R$ free, giving $c_{\mathrm{neg}} + c_{\mathrm{div}} + c_{\exp} + c_{\mathrm{mul}}
= 2 + 1 + 1 + 2 = 6 - 1 = 5$.
\end{example}

% ── Example 5 ───────────────────────────────────────────────────────────────
\begin{example}[Von Bertalanffy Growth --- deep nesting, no penalty]
\label{ex:vonb}
\[
  L(t) \;=\; L_\infty \cdot \bigl(1 - \exp(-K \cdot (t - t_0))\bigr),
\]
where $L_\infty$, $K$, $t_0$ are constants and $t$ is the free variable.

\textbf{Operation tree.}
\begin{align*}
  & \texttt{sub}(t,\, t_0)                  &\text{[cost 2]} \\
  & \texttt{mul}(K,\, \cdots)               &\text{[cost 2]} \\
  & \texttt{neg}(\cdots)                     &\text{[cost 2]} \\
  & \texttt{exp}(\cdots)                     &\text{[cost 1]} \\
  & \texttt{sub}(1,\, \cdots)               &\text{[cost 2]} \\
  & \texttt{mul}(L_\infty,\, \cdots)        &\text{[cost 2]}
\end{align*}

\textbf{Na\"{i}ve cost.}
\[
  \NC(L) \;=\; 2 + 2 + 2 + 1 + 2 + 2 \;=\; 11.
\]

\textbf{Sharing discount.}
No subexpression appears at multiple sites; $K$ and $t_0$ are free
constants, not computed intermediates.  $\SD(L) = 0$.

\textbf{Pattern bonus.}
No compound $\mathcal{F}_{16}$ pattern matches a subexpression of $L(t)$.
$\PB(L) = 0$.

\textbf{Decomposition.}
\[
  \Cost(L) \;=\; 11 - 0 - 0 \;=\; 11. \qquad \checkmark
\]

\textbf{No nesting penalty.}
The expression has nesting depth 6 (six nested operators from
$\texttt{mul}(L_\infty, \ldots)$ down to $\texttt{sub}(t, t_0)$).
By Lemma~\ref{lem:nesting} and Corollary~\ref{cor:additivity}, the cost
is simply the sum of the six unit costs: $2+2+2+1+2+2 = 11$.
Depth contributes no overhead beyond the sum of parts.
\end{example}

% ════════════════════════════════════════════════════════════════════════════
\section{Summary of Worked Examples}
\label{sec:summary-table}
% ════════════════════════════════════════════════════════════════════════════

\begin{table}[h]
\centering
\caption{T38 decomposition verified on five catalogue expressions.}
\label{tab:examples}
\smallskip
\begin{tabular}{@{}lccccl@{}}
\toprule
\textbf{Expression} & $\NC$ & $\SD$ & $\PB$ & $\Cost$ & \textbf{Key feature} \\
\midrule
Ohm's Law ($V=IR$)                 &  2 & 0 & 0 &  2 & tight bound \\
Allometric ($Y=aM^b$)              &  5 & 0 & 0 &  5 & tight, nested \\
pH ($-\log_{10}[\mathrm{H}^+]$)   &  5 & 2 & 0 &  3 & constant folding \\
Arrhenius ($Ae^{-E_a/RT}$)         &  6 & 1 & 0 &  5 & constant folding \\
Von Bertalanffy ($L_\infty(1-e^{-K(t-t_0)})$) & 11 & 0 & 0 & 11 & depth-6 nesting \\
\bottomrule
\end{tabular}
\end{table}

In all five cases the formula $\Cost = \NC - \SD - \PB$ gives the exact
correct result with non-negative, precisely computable discount terms.
The Von Bertalanffy example with depth-6 nesting explicitly confirms the
No Nesting Penalty (Lemma~\ref{lem:nesting}).

% ════════════════════════════════════════════════════════════════════════════
\section{Theorem Catalog Labels}
\label{sec:catalog}
% ════════════════════════════════════════════════════════════════════════════

\begin{description}
  \item[T33] \emph{Sub 2-Node Optimality.}
        $\Cost(\mathrm{sub}) = 2$ in $\mathcal{F}_{16}$, achieved by
        $\LEdiv(x, \EML(y,1)) = \ln(e^x/e^y) = x-y$.
        Reference: \textit{T33\_Sub\_2Node.tex}.

  \item[T34] \emph{Na\"{i}ve Upper Bound.}
        $\Cost(E) \leq \NC(E)$ for all EML-computable $E$.
        Reference: COST-2, \textit{COST2\_Naive\_Upper\_Bound.tex}.

  \item[T35--T37] \emph{Lower Bounds} (trivial, depth, operation-type).
        Reference: COST-3, \textit{COST3\_Lower\_Bound.tex}.

  \item[T38] \emph{Cost Decomposition.}
        $\Cost(E) = \NC(E) - \SD(E) - \PB(E)$, with
        $\SD(E) \geq 0$ (sharing discount, COST-5) and
        $\PB(E) \geq 0$ (pattern bonus, COST-7).
        No Nesting Penalty: $\Cost(\mathrm{op}_1(\mathrm{op}_2(x)))
        = c_{\mathrm{op}_1} + c_{\mathrm{op}_2}$.
        Reference: this document (COST-6).

  \item[T38-NNP] \emph{No Nesting Penalty} (Lemma~\ref{lem:nesting}).
        EML operator composition is free: nesting depth adds no
        representation cost beyond the sum of parts.
        Reference: this document (COST-6, Section~\ref{sec:nesting}).
\end{description}

% ════════════════════════════════════════════════════════════════════════════
\section{Roadmap}
\label{sec:roadmap}
% ════════════════════════════════════════════════════════════════════════════

\begin{center}
\begin{tabular}{@{}ll@{}}
\toprule
\textbf{Session} & \textbf{Topic} \\
\midrule
COST-1 & Cost model foundations (SuperBEST v3 unit costs) \\
COST-2 & Na\"{i}ve upper bound $\NC(E)$ --- Theorem T34 \\
COST-3 & Lower bounds --- Theorems T35--T37 \\
COST-4 & Constant-folding discount rules \\
COST-5 & Sharing discount $\SD(E)$ --- formal theory \\
COST-6 & \textbf{Cost Decomposition --- Theorem T38 (this document)} \\
COST-7 & Pattern bonus $\PB(E)$ --- catalogue and algorithm \\
COST-8 & Blind test: 20-equation validation (\textit{COST8\_Blind\_Test.md}) \\
COST-9 & General cost optimisation algorithm \\
\bottomrule
\end{tabular}
\end{center}

% ════════════════════════════════════════════════════════════════════════════
\section*{Acknowledgments}
% ════════════════════════════════════════════════════════════════════════════

This work is part of the monogate SuperBEST cost theory sprint (2026-04-20).
Theorem T38 synthesises the upper bound (T34, COST-2) and the reduction
framework (sharing discount, COST-5; pattern bonus, COST-7) into a single
exact equality.  The No Nesting Penalty (T38-NNP) resolves a natural
question about the overhead of operator composition in the EML family and
is confirmed by the COST-8 blind test (session 16, Von Bertalanffy).

\bigskip
\noindent\textit{Almaguer, A.R.\ (2026).
Cost Theory COST-6: The Cost Decomposition Theorem (T38).
monogate research. \url{https://monogate.org}}

\end{document}
